% Part: lambda-calculus
% Chapter: introduction
% Section: unique-readability

\documentclass[../../../include/open-logic-section]{subfiles}

\begin{document}

\olfileid{lam}{syn}{unq}
\olsection{اکّو طرح پڑھیا جانا}

اسیں پُچھ سکدے آں کہ کی ہر ٹرم بناؤن دا اکّو طریقہ اے۔ جواب ہاں
اے۔ ہر لامبڈا ٹرم بناؤن تے سمجھّن دا اکّو طریقہ اے۔ اگلی گل بات وچ
\emph{بناوٹ} توں مراد \olref[trm]{defn:term} دے قاعدے اک یا کئی
وار ورت کے ٹرم بناؤن دی کاروائی اے۔

\begin{lem}\ollabel{lem:term-start}
ٹرم متغیر یا قوس نال شروع ہوندا اے۔
\end{lem}

\begin{proof}
کوئی شے ٹرم اُدوں ای منّی جاندی اے جدوں اوہ \olref[trm]{defn:term}
مطابق بنی ہووے۔ جے ایہ \olref[trm]{defn:term-var} دا نتیجہ اے، تاں
لازماً متغیر اے۔ جے ایہ \olref[trm]{defn:term-abs} یا
\olref[trm]{defn:term-app} دا نتیجہ اے، تاں قوس نال شروع ہوندا اے۔
\end{proof}

\begin{lem}\ollabel{lem:app-start}
اطلاق دا نتیجہ یا تاں دو قوساں نال، یا اک قوس تے اک متغیر نال
شروع ہوندا اے۔
\end{lem}

\begin{proof}
جے $M$ اطلاق دا نتیجہ اے، تاں اوہدی شکل $(PQ)$ اے، سو اوہ قوس نال
شروع ہوندا اے۔ کیوں جے $P$ اک ٹرم اے، ایس لئی \olref{lem:term-start}
مطابق اوہ قوس یا متغیر نال شروع ہوندا اے۔
\end{proof}

\begin{lem}\ollabel{lem:initial}
کسے ٹرم دا کوئی اوہ مُڈھلا حصہ جیہڑا پورے ٹرم توں چھوٹا ہووے،
آپ ٹرم نہیں ہو سکدا۔
\end{lem}

\begin{prob}
ٹرماں دی لمبائی اُتے استقرا نال \olref[lam][syn][unq]{lem:initial}
ثابت کرو۔
\end{prob}

\begin{prop}[اکّو طرح پڑھیا جانا] \ollabel{prop:unq}
ہر ٹرم دی بناوٹ اکّو اے۔ دوجے لفظاں وچ، جے ٹرم $M$ کسے بناوٹ نال
بنیا اے، تاں اوہو اکّو بناوٹ اے جیہڑی ایہ ٹرم بنا سکدی اے۔
\end{prop}

\begin{proof}
اسیں ٹرماں دی بناوٹ اُتے استقرا نال ثابت کردے آں۔
\begin{enumerate}
  \item $M$ دی شکل $x$ اے، جتھے $x$ کوئی متغیر اے۔ کیوں جے
    تجرید تے اطلاق دے نتیجے ہمیشہ قوساں نال شروع ہوندے نیں،
    اوہناں نال $M$ نہیں بن سکدا۔ سو $M$ دی بناوٹ لازماً
    \olref[trm]{defn:term}\olref[trm]{defn:term-var} دا اکّو قدم اے۔
  \item $M$ دی شکل $(\lambd[x][N])$ اے، جتھے $x$ کوئی متغیر
    تے $N$ اک ٹرم اے۔ ایہ \olref[trm]{defn:term}\olref[trm]{defn:term-var}
    نال نہیں بن سکدا، کیوں جے ایہ اکّلا متغیر نہیں۔
    \olref{lem:app-start} مطابق ایہ اطلاق دا نتیجہ وی نہیں۔
    سو $M$ صرف~$N$ دی تجرید نال بنیا ہو سکدا اے۔ استقرا دے مفروضے
    توں پتا اے کہ $N$ دی اپنی بناوٹ وی اکّو اے۔
  \item $M$ دی شکل $(PQ)$ اے، جتھے $P$ تے $Q$ ٹرم نیں۔
    کیوں جے ایہ قوس نال شروع ہوندا اے، ایس لئی
    \olref[trm]{defn:term}\olref[trm]{defn:term-var} نال نہیں
    بن سکدا۔ \olref{lem:term-start} مطابق $P$، $\lambd$ نال
    شروع نہیں ہو سکدا، سو $(PQ)$ تجرید دا نتیجہ نہیں ہو سکدا۔
    ہُن منّ لو $M$ اطلاق نال کسے ہور طریقے وی بن سکدا اے،
    مثلاً ایہ $(P'Q')$ دی شکل دا وی اے۔ فیر $P$، $P'$ دا
    پورے توں چھوٹا مُڈھلا حصہ اے، یا ایہدا اُلٹ اے۔ پر
    \olref{lem:initial} مطابق ایہ ممکن نہیں۔ سو $P$ تے $Q$
    اکّو طرح متعین نیں، تے استقرا دے مفروضے توں پتا اے کہ
    $P$ تے $Q$ دی بناوٹ وی اکّو اے۔
\end{enumerate}
\end{proof}

اُپرلے قضیے نوں سوکھے لفظاں وچ ایویں آکھ سکدے آں:
\begin{prop}
  ٹرم $M$ دی شکل صرف ہِٹھلیاں وچوں اک ہو سکدی اے:
  \begin{enumerate}
    \item $x$، جتھے $x$ اوہ متغیر اے جیہڑا $M$ توں اکّو طرح متعین اے۔
    \item $(\lambd[x][N])$، جتھے $x$ متغیر اے تے $N$ دوجا ٹرم۔
      دوویں $M$ توں اکّو طرح متعین نیں۔
    \item $(PQ)$، جتھے $P$ تے $Q$ دو ٹرم نیں جیہڑے $M$ توں
      اکّو طرح متعین نیں۔
  \end{enumerate}
\end{prop}

\end{document}
